% ---------------------------------------------------------------------------
% Conditional recourse: local filtering, its limits, exact cut recourse,
% core search, and exact output.  Self-contained fragment (no preamble).
% Shared notation: X = prod X_i, X_i = [l_i,u_i] or [l_i,u_i] cap Z; F, \OPT=F^*;
% L = upper coordinate curvature; g = growth; kappa = max{1,L/g}; p = bag size.
% Citation keys not yet in references.bib: Hammer1965, PicardRatliff1975,
% Topkis1978, Edmonds1970, BeierVocking2004, RoglinVocking2007.
% ---------------------------------------------------------------------------

\subsection{Conditional values and bag-cell bounds}\label{sec:cr-values}

Throughout this part $B\subseteq\{1,\ldots,n\}$ is a nonempty set of
coordinates, for example a bag of the tree decomposition or a supplied
core, and $-B$ is its complement. Write
$\widehat X_B=\prod_{i\in B}[l_i,u_i]$ for the continuous hull of $X_B$.
As in the rest of the paper, $F$ is defined on the continuous hull of $X$
and $L$ is a valid upper coordinate curvature there.

\begin{definition}[Conditional value]\label{def:cr-W}
For $v\in\widehat X_B$ put
\[
W_B(v)=\min\{F(v,x_{-B}) : x_{-B}\in X_{-B}\}.
\]
The minimum exists because $F$ is continuous and $X_{-B}$ is compact.
Clearly $W_B(x_B)\le F(x)$ for $x\in X$, and $\min_{v\in X_B}W_B(v)=\OPT$.
\end{definition}

A \emph{bag cell} is a product $C=\prod_{i\in B}[a_i,b_i]$ with
$a_i\le b_i$ and $a_i,b_i\in X_i$. Its \emph{corners} are the points $v$
with $v_i\in\{a_i,b_i\}$; they are feasible bag values. Put
\[
e_B(C)=\frac L8\sum_{i\in B}w_i(C)^2,\qquad
w_i(C)=\begin{cases}
0,& i\in I_Z\text{ and }b_i-a_i=1,\\
b_i-a_i,&\text{otherwise.}
\end{cases}
\]
If every side of $C$ is at most $h$, then $e_B(C)\le|B|Lh^2/8$.

\begin{lemma}[Conditional values inherit coordinate curvature]\label{lem:cr-semiconcave}
For every $i\in B$ and every fixed $v_{B\setminus\{i\}}$, the function
$t\mapsto W_B(v_{B\setminus\{i\}},t)-Lt^2/2$ is concave on $[l_i,u_i]$.
\end{lemma}
\begin{proof}
Fix $x_{-B}\in X_{-B}$. The point $(v_{B\setminus\{i\}},t,x_{-B})$ lies in
the continuous hull of $X$ for $t\in[l_i,u_i]$, so
$t\mapsto F(v_{B\setminus\{i\}},t,x_{-B})-Lt^2/2$ is concave. The
function in the statement is the pointwise minimum of these concave
functions over the set $X_{-B}$, which does not depend on $t$. A pointwise
infimum of concave functions is concave, because its hypograph is the
intersection of their hypographs, and here it is finite.
\end{proof}

\begin{lemma}[Bag-cell interpolation]\label{lem:cr-cell}
Let $C$ be a bag cell. Every $x\in X$ with $x_B\in C$ satisfies
\[
F(x)\ \ge\ \min_{v\text{ a corner of }C}W_B(v)-e_B(C).
\]
\end{lemma}
\begin{proof}
Put $\xi=x_B$. Let $(Y_i)_{i\in B}$ be independent, with $Y_i\in\{a_i,b_i\}$
and $\E Y_i=\xi_i$ (and $Y_i=a_i$ if $a_i=b_i$). Then
$\operatorname{Var}Y_i=(\xi_i-a_i)(b_i-\xi_i)\le(b_i-a_i)^2/4$. If $i$ is an
integer coordinate with $b_i-a_i=1$, then $\xi_i\in\{a_i,b_i\}$ because
$x$ is feasible, and $\operatorname{Var}Y_i=0$. Hence
$\frac L2\sum_i\operatorname{Var}Y_i\le e_B(C)$.

Order $B$ as $i_1,\ldots,i_r$, and let $Z^{(s)}$ be $\xi$ with its first $s$
listed coordinates replaced by the corresponding $Y$'s. All $Z^{(s)}$ lie
in $\widehat X_B$. Conditionally on $Y_{i_1},\ldots,Y_{i_{s-1}}$, the
function $\varphi(t)=W_B(Z^{(s-1)}\text{ with coordinate }i_s\text{ set to }t)$
has $\varphi(t)-Lt^2/2$ concave by Lemma~\ref{lem:cr-semiconcave}. Jensen's
inequality for this concave function gives
$\E\varphi(Y_{i_s})\le\varphi(\xi_{i_s})+\frac L2\operatorname{Var}Y_{i_s}$,
that is,
$\E[W_B(Z^{(s)})\mid Y_{i_1},\ldots,Y_{i_{s-1}}]
\le W_B(Z^{(s-1)})+\frac L2\operatorname{Var}Y_{i_s}$.
Taking expectations and summing over $s$ yields
$\E W_B(Y)\le W_B(\xi)+e_B(C)$. Since $Y$ is a corner of $C$ almost surely,
\[
\min_{\text{corners}}W_B\le\E W_B(Y)\le W_B(\xi)+e_B(C)\le F(x)+e_B(C).
\qedhere
\]
\end{proof}

No coordinate outside $B$ is rounded, so the error $e_B(C)$ does not grow
with $n$. The price is that $W_B$ is a global optimization problem over
the original outside domain.

\subsection{The conditional-recourse filter}\label{sec:cr-interface}

\begin{theorem}[Conditional-recourse filter]\label{thm:cr-filter}
Let $\mathcal Q$ be a finite family of bag cells and
$e\ge\max_{C\in\mathcal Q}e_B(C)$. Suppose that for every corner $v$ of
every cell of $\mathcal Q$ an oracle supplies a number $\ell(v)$ and a
point $x^v\in X$ with $x^v_B=v$ such that
\begin{equation}\label{eq:cr-oracle}
\ell(v)\le W_B(v),\qquad F(x^v)\le\ell(v)+\eta .
\end{equation}
Let $U$ be the value of a feasible point with $U\le\min_vF(x^v)$, the
minimum taken over all supplied completions. Retain $C\in\mathcal Q$ if
and only if $\min_{v\text{ corner of }C}\ell(v)-e\le U$. Then:
\begin{enumerate}
\item[(a)] every $C\in\mathcal Q$ that contains $x^*_B$ for some global
optimizer $x^*$ is retained;
\item[(b)] if some $C\in\mathcal Q$ contains $x^*_B$ for a global optimizer
$x^*$, then $U\le\OPT+e+\eta$, and every retained cell has a corner $v$
with $W_B(v)\le\OPT+2e+2\eta$;
\item[(c)] every $x\in X$ with $x_B\in\bigcup\mathcal Q$ satisfies
\[
F(x)\ge\lambda:=\min_{C\in\mathcal Q}\Bigl(\min_{\text{corners of }C}\ell-e\Bigr),
\]
and $U-\lambda\le e+\eta$.
\end{enumerate}
\end{theorem}
\begin{proof}
(a) By~\eqref{eq:cr-oracle} and Lemma~\ref{lem:cr-cell} applied to $x^*$,
\[
\min_{\text{corners}}\ell-e\le\min_{\text{corners}}W_B-e_B(C)\le\OPT\le U.
\]

(b) Lemma~\ref{lem:cr-cell} at $x^*$ gives a corner $v^\circ$ of the
optimizer's cell with $W_B(v^\circ)\le\OPT+e$. Hence
$U\le F(x^{v^\circ})\le\ell(v^\circ)+\eta\le W_B(v^\circ)+\eta\le\OPT+e+\eta$.
A retained cell has a corner $v$ with $\ell(v)\le U+e$, and
$W_B(v)\le F(x^v)\le\ell(v)+\eta\le U+e+\eta\le\OPT+2e+2\eta$.

(c) For $x_B\in C$, Lemma~\ref{lem:cr-cell} and $\ell\le W_B$ give
$F(x)\ge\min_{\text{corners}}W_B-e\ge\min_{\text{corners}}\ell-e\ge\lambda$.
Every corner satisfies $\ell(v)\ge F(x^v)-\eta\ge U-\eta$, so
$\lambda\ge U-e-\eta$.
\end{proof}

\begin{remark}
Hypothesis~\eqref{eq:cr-oracle} is a certified global optimization of the
outside problem, over the original outside domain, at every queried bag
value, with an error $\eta$ that is uniform over these values. In a
multilevel scheme that refines only retained cells, the hypothesis of
(b) holds at every level by (a) and induction. A point whose bag
projection lies in a cell discarded at an earlier level has value larger
than the incumbent at that level, hence larger than the current $U$.
Therefore $\min\{\lambda,U\}$ is a certified lower bound on $\OPT$, as in
the filtering proposition for corrected grids.
\end{remark}

The next proposition shows that the filter needs only control of the
\emph{variation} of the error across bag states. A common additive error,
such as a normalization constant of messages, is harmless.

\begin{proposition}[Only the error oscillation matters]\label{prop:cr-osc}
Let $\mathcal Q$, $e$ be as in Theorem~\ref{thm:cr-filter}, and suppose
some cell of $\mathcal Q$ contains $x^*_B$ for a global optimizer $x^*$.
Let $\mathcal V$ be the set of all corners and $\ell:\mathcal V\to\R$ a
function such that, for some constants $\delta_-\le\delta_+$ with
$\delta_+-\delta_-\le\eta$,
\[
\delta_-\le W_B(v)-\ell(v)\le\delta_+\qquad(v\in\mathcal V).
\]
The constants need not be known. Retain $C$ if and only if
$\min_{\text{corners of }C}\ell\le\min_{\mathcal V}\ell+e+\eta$. Then every
cell containing $x^*_B$ for a global optimizer $x^*$ is retained, and
every retained cell has a corner with $W_B(v)\le\OPT+2e+2\eta$.
\end{proposition}
\begin{proof}
Let $C^*$ contain $x^*_B$. By Lemma~\ref{lem:cr-cell},
\[
\min_{\text{corners of }C^*}\ell\le\min_{\text{corners of }C^*}W_B-\delta_-
\le\OPT+e-\delta_-,
\]
while $W_B\ge\OPT$ gives
$\min_{\mathcal V}\ell\ge\OPT-\delta_+$. The difference is at most
$e+\delta_+-\delta_-\le e+\eta$, so $C^*$ is retained. If $C$ is retained,
it has a corner $v$ with
$W_B(v)\le\ell(v)+\delta_+\le\min_{\mathcal V}\ell+e+\eta+\delta_+$, and
$\min_{\mathcal V}\ell\le\min_{\text{corners of }C^*}\ell\le\OPT+e-\delta_-$.
Thus $W_B(v)\le\OPT+2e+\eta+(\delta_+-\delta_-)\le\OPT+2e+2\eta$.
\end{proof}

Theorem~\ref{thm:cr-filter} additionally certifies a gap; for that, valid
lower bounds and feasible completions are needed. For filtering alone,
Proposition~\ref{prop:cr-osc} identifies the essential requirement: the
oscillation of $W_B-\ell$ over the compared bag states must be
$O(e)$.

\subsection{Grid min-marginals require a dimension-dependent budget}\label{sec:cr-star}

Let $h=2^{-j}$ and $G=h\Z\cap[0,1]$. On $[0,1]^{1+m}$ with coordinates
$(x,y_1,\ldots,y_m)$, let the uniform product grid be $G^{1+m}$, write
$U_G=\min_{G^{1+m}}F$, and, for the bag $B=\{x,y_1\}$ and $v\in G^2$,
\[
m^F(v)=\min\{F(z):z\in G^{1+m},\ (z_x,z_{y_1})=v\}.
\]
A \emph{constant-correction test} with constant $c$ retains a bag cell
$C$ of the grid if and only if $\min_{\text{corners of }C}m^F-c\le U_G$;
its coordinate version retains a grid interval $[a,a+h]$ of the centre
coordinate if and only if $\min\{V_h(a),V_h(a+h)\}-c\le U_G$, where
$V_h(a)=\min\{F(z):z\in G^{1+m},\ z_x=a\}$. On a uniform grid every node has
correction $Lh^2/8$ in every coordinate, so the corrected-grid method of
the paper applies the coordinate version with $c=(1+m)Lh^2/8$. A bag-local
budget would take $c=|B|Lh^2/8=Lh^2/4$. In both families below, the
corner minimum of each relevant bag cell equals
$\min\{V_h(a),V_h(a+h)\}$, where $[a,a+h]$ is the projection of the cell
on the centre coordinate. Hence the thresholds stated for bag cells hold
verbatim for the coordinate version.

\begin{proposition}[Constant corrections must grow with dimension]\label{prop:cr-star}
\leavevmode
\begin{enumerate}
\item[(A)] Let $j\ge1$, $m\ge1$, $\epsilon>0$ with $mh^2>4(1+\epsilon)$, and
\[
F_A(x,y)=x^2-hx\sum_{i=1}^my_i+\sum_{i=1}^my_i^2
+\Bigl(\frac{mh^2}4-1-\epsilon\Bigr)x .
\]
Its interaction graph is a star with bags $\{x,y_i\}$ ($p=2$), and $L=2$
is valid. Its unique minimizer is $x=1$, $y_i=h/2$, with $\OPT=-\epsilon$,
and any growth constant satisfies $g\le\epsilon/(1+mh^2/4)$; moreover
$U_G=0$. The only bag cell containing $(1,h/2)$ is
$C=[1-h,1]\times[0,h]$, and
\[
\min_{\text{corners of }C}m^F=(1-h)\Bigl(\frac{mh^2}4-h-\epsilon\Bigr).
\]
Hence a constant-correction test discards a cell containing the unique
optimizer unless
$c\ge(1-h)\bigl(m-4/h-4\epsilon/h^2\bigr)\,Lh^2/8$. As $m\to\infty$ the
ratio of $(1+m)Lh^2/8$ to this lower bound tends to $(1-h)^{-1}$.
\item[(B)] Let $j\ge2$, let $d\ge2/h$ be an integer, $m=d^2$, and
\[
F_B(x,y)=\Bigl(x-\frac12\Bigr)^2+\sum_{i=1}^m
\Bigl(y_i-\frac h2-\frac{x-1/2}{d}\Bigr)^2 .
\]
It is a star with $p=2$; $L=4$ is valid; its unique minimizer
$z^*=(1/2,h/2,\ldots,h/2)$ is interior; and
$F_B(z)-\OPT\ge\frac{3-\sqrt5}2\norm{z-z^*}^2$, so
$\kappa\le2(3+\sqrt5)<11$ for every $d$ and $h$. Moreover
$U_G=\frac12-\frac{dh}2+\frac{d^2h^2}4$, and both bag cells containing
$(1/2,h/2)$ have
\[
\min_{\text{corners}}m^F-U_G=dh\Bigl(\frac12-h\Bigr)+2h^2-\frac12 .
\]
Thus a constant-correction test discards every cell containing the
unique optimizer unless $c\ge dh(\frac12-h)+2h^2-\frac12$, a quantity of
order $\sqrt m\,h$ at fixed $h$, while $p=2$ and $\kappa<11$.
\end{enumerate}
\end{proposition}
\begin{proof}
(A) Write $\beta=mh^2/4-1-\epsilon>0$. All coordinate second derivatives
equal $2$ and $F_A$ is quadratic, so $L=2$ is valid. For fixed $x$, each
leaf term $y^2-hxy$ is strictly convex with minimizer $hx/2\in[0,1]$ and
minimum $-h^2x^2/4$. Hence $\min_yF_A(x,y)=(1-mh^2/4)x^2+\beta x=:V(x)$,
which is strictly concave because $mh^2>4$, with $V(0)=0$ and
$V(1)=-\epsilon$. For $x\in(0,1)$, strict concavity gives
$V(x)>(1-x)V(0)+xV(1)\ge-\epsilon$. So $x=1$ is the unique minimizing
centre and $y_i=h/2$ the unique leaves; $\OPT=-\epsilon$. The origin has
$F_A=0=\OPT+\epsilon$ and squared distance $1+mh^2/4$ from the minimizer,
which bounds $g$.

On the grid, $y(y-hx)\ge0$ for $y\in G$ and $x\in[0,1]$, since either
$y=0$ or $y\ge h\ge hx$. Also $x^2+\beta x\ge0$ for $x\ge0$. Hence
$F_A\ge0$ on $G^{1+m}$ with equality at the origin, so $U_G=0$, and every
leaf's grid minimum is $0$. The point $(1,h/2)$ has $x$ at the right
endpoint and $h/2$ interior to $[0,h]$, so $C$ is the only grid cell of
the bag containing it. At its corners,
$m^F(x,y_1)=x^2+\beta x+(y_1^2-hxy_1)$. For $x=1-h$ this equals
$(1-h)(1-h+\beta)=(1-h)(mh^2/4-h-\epsilon)$ at $y_1=0$ and exceeds it by
$h^3$ at $y_1=h$. For $x=1$ it equals $mh^2/4-\epsilon$ at both $y_1=0,h$,
which exceeds the previous value by $h(mh^2/4-\epsilon)+h(1-h)>0$. This
proves the corner minimum; it also equals $\min\{V_h(1-h),V_h(1)\}$,
because every leaf's grid minimum is attained at $0$. Since $U_G=0$ and
$L=2$, the threshold for $c$ is as stated, and
$(1+m)h^2/4\big/\bigl[(1-h)(mh^2/4-h-\epsilon)\bigr]\to(1-h)^{-1}$.

(B) The Hessian $H$ has $H_{xx}=2+2m/d^2=4$, $H_{y_iy_i}=2$,
$H_{xy_i}=-2/d$, and zero elsewhere; thus $L=4$ is valid. Vectors with
zero $x$-component and leaf components summing to zero are eigenvectors
with eigenvalue $2$. On the span of $e_x$ and $m^{-1/2}\mathbf 1_y$, $H$
acts as $\bigl(\begin{smallmatrix}4&-2\\-2&2\end{smallmatrix}\bigr)$, with
eigenvalues $3\pm\sqrt5$. Since $F_B\ge0$ vanishes only at $z^*\in(0,1)^{1+m}$
and $F_B(z)=\frac12(z-z^*)^{\mathsf T}H(z-z^*)$, growth holds with
$g=(3-\sqrt5)/2$ and $\kappa=L/g=2(3+\sqrt5)$.

Write $t=x-1/2\in h\Z\cap[-1/2,1/2]$ for centre grid values (here
$1/2\in G$ because $j\ge1$). A leaf's ideal value
$\iota(t)=h/2+t/d$ lies in $[0,h]$ because $|t|/d\le1/(2d)\le h/4$.
Its distance to $G$ is therefore $h/2-|t|/d$, and the conditional grid
value of the centre is
\[
V_h(t)=t^2+m\Bigl(\frac h2-\frac{|t|}d\Bigr)^2=2t^2-dh|t|+\frac{d^2h^2}4 .
\]
On the grid $|t|\le1/2\le dh/4$, and $s\mapsto2s^2-dhs$ decreases on
$[0,dh/4]$, so $U_G=V_h(\pm1/2)=\frac12-\frac{dh}2+\frac{d^2h^2}4$.
The optimizer's bag projection $(1/2,h/2)$ lies on the common edge of
the cells $[1/2-h,1/2]\times[0,h]$ and $[1/2,1/2+h]\times[0,h]$ and in no
other cell. At $t=0$ both leaf labels $0,h$ are at distance $h/2$ from
$\iota$, so $m^F=d^2h^2/4$. At $t=\pm h$ the nearer label gives
$m^F=V_h(h)=2h^2-dh^2+d^2h^2/4$ and the farther adds $2h^2/d$. Since
$d\ge2$, each cell has corner minimum $d^2h^2/4-(d-2)h^2
=\min\{V_h(0),V_h(\pm h)\}$, and subtracting $U_G$ gives the stated
difference. For $h\le1/4$ this difference is at least
$dh/4+2h^2-1/2$, which grows linearly in $d=\sqrt m$. Discarding both
centre intervals adjacent to $x=1/2$ removes the optimizer.
\end{proof}

\begin{example}[The smallest instance]\label{ex:cr-star32}
Take $h=1/2$, $m=32$, $\epsilon=1/16$ in (A):
$F=x^2-\frac x2\sum y_i+\sum y_i^2+\frac{15}{16}x$ on $[0,1]^{33}$.
The exact leaf recourse is $y_i=x/4$, with $V(x)=-x^2+\frac{15}{16}x$, so
the unique optimizer is $x=1$, $y_i=1/4$, value $-1/16$. On
$G=\{0,\frac12,1\}$ the conditional grid values at $x=0,\frac12,1$ are
$0,\frac{23}{32},\frac{31}{16}$, and $U_G=0$. The four corner
min-marginals of $C=[\frac12,1]\times[0,\frac12]$ are
$\frac{23}{32},\frac{27}{32},\frac{31}{16},\frac{31}{16}$. The bag-local
budget $\frac18$ rejects $C$; the global correction $\frac{33}{16}$ keeps it.
The conditional grid error is
\[
V_h(x)-V(x)=m\,\dist(x/4,G)^2,
\]
which varies from $0$ at $x=0$ to $m/16$ at $x=1$; no normalization of
messages removes this variation. Because the growth constant of family (A)
degrades as $m$ grows, this instance refutes budgets depending on
$(p,L)$ only; family (B) refutes budgets depending on $(p,\kappa)$.
\end{example}

\subsection{An exact minimum-cut oracle for a large nonconvex residual}\label{sec:cr-cut}

In this subsection
\[
F(x)=c_0+\sum_i(a_ix_i^2+b_ix_i)+\sum_{i<j}c_{ij}x_ix_j
\]
is a rational quadratic. Its coordinates are split into a \emph{core}
$\mathcal C$ of $k$ continuous coordinates and a \emph{residual}
$\mathcal R$. Missing couplings have $c_{ij}=0$. Assume
\begin{equation}\label{eq:cr-class}
\text{(R1)}\ a_i\le0\ (i\in\mathcal R),\qquad
\text{(R2)}\ c_{ij}o_io_j\le0\ (i,j\in\mathcal R)\ \text{for some }
o\in\{-1,1\}^{\mathcal R}.
\end{equation}
Nothing is assumed about core--core or core--residual coefficients, and
the residual interaction graph may be dense. A residual coordinate may be
continuous or integer; integer bounds are rounded inward first and an
empty domain is rejected. Given the signs of the residual $c_{ij}$, a
breadth-first traversal finds an orientation satisfying (R2) or a cycle
proving that none exists.

\begin{lemma}[Endpoint reduction]\label{lem:cr-endpoint}
Fix $v\in\widehat X_{\mathcal C}$ and a box $\prod_{i\in\mathcal R}[l_i,u_i]$
(after inward rounding). Under (R1), the minimum of $F(v,\cdot)$ over the
continuous box is attained at a point with $x_i\in\{l_i,u_i\}$ for all
$i\in\mathcal R$. Consequently it equals the minimum over the mixed
residual domain.
\end{lemma}
\begin{proof}
Start from a minimizer, which exists by compactness. For each residual
coordinate in turn, the objective as a function of that coordinate alone
is $a_it^2+(\text{affine in }t)$, which is concave by (R1); replace the
coordinate by an endpoint with no larger value. The result is an endpoint
minimizer. Its coordinates are integral for integer coordinates, so it is
feasible for the mixed domain, which is contained in the continuous box.
\end{proof}

For $i\in\mathcal R$ put $\alpha_i=l_i$, $d_i=u_i-l_i$ if $o_i=1$, and
$\alpha_i=u_i$, $d_i=-(u_i-l_i)$ if $o_i=-1$. For $y\in\{0,1\}^{\mathcal R}$
let $x_i(y)=\alpha_i+d_iy_i$; this enumerates all residual endpoint
vectors. In the network below, for a node set $S$ that contains the
source $\mathsf s$ but not the sink $\mathsf t$, $\operatorname{cap}(S)$
denotes the total capacity of arcs leaving $S$.

\begin{proposition}[Signed cut representation]\label{prop:cr-cut}
Fix a rational core point $v$. Under~\eqref{eq:cr-class} there are rationals
$\varphi_0,\mu_i,\omega_{ij}$, computable by exact evaluation of $F$, such
that for all $y\in\{0,1\}^{\mathcal R}$
\[
F(v,x(y))=\varphi_0+\sum_{i}\mu_iy_i+\sum_{i<j}\omega_{ij}y_iy_j,
\qquad \omega_{ij}=c_{ij}d_id_j\le0 .
\]
Let $\rho_i=\mu_i+\frac12\sum_{j\ne i}\omega_{ij}$ and
$r_{ij}=-\omega_{ij}/2\ge0$. Build a network on
$\mathcal R\cup\{\mathsf s,\mathsf t\}$ with arcs $i\to\mathsf t$ of
capacity $\max\{\rho_i,0\}$, arcs $\mathsf s\to i$ of capacity
$\max\{-\rho_i,0\}$, and, for each pair with $\omega_{ij}\ne0$, arcs
$i\to j$ and $j\to i$ of capacity $r_{ij}$. With
$S_y=\{\mathsf s\}\cup\{i:y_i=1\}$,
\begin{equation}\label{eq:cr-cutid}
F(v,x(y))=\varphi_0+\sum_i\min\{0,\rho_i\}+\operatorname{cap}(S_y).
\end{equation}
\end{proposition}
\begin{proof}
Expand $F(v,\alpha+d\circ y)$. Core terms and $c_0$ give constants. A term
$a_ix_i^2$ gives $a_i\alpha_i^2+(2a_i\alpha_id_i+a_id_i^2)y_i$ by
$y_i^2=y_i$; $b_ix_i$ and core--residual terms $c_{ij}v_jx_i$ are affine
in $y_i$; and $c_{ij}x_ix_j$ for residual $i,j$ gives
$c_{ij}d_id_jy_iy_j$ plus affine terms. By (R2),
$\omega_{ij}=c_{ij}o_io_j(u_i-l_i)(u_j-l_j)\le0$. For binary $y$,
$\omega y_iy_j=\frac\omega2(y_i+y_j)-\frac\omega2|y_i-y_j|$, so
\[
F(v,x(y))=\varphi_0+\sum_i\rho_iy_i+\sum_{i<j}r_{ij}|y_i-y_j| .
\]
In the cut $S_y$, an arc $i\to\mathsf t$ is cut exactly when $y_i=1$, an
arc $\mathsf s\to i$ exactly when $y_i=0$, and exactly one of the two pair
arcs exactly when $y_i\ne y_j$. Hence
$\operatorname{cap}(S_y)=\sum_i[\max\{\rho_i,0\}y_i+\max\{-\rho_i,0\}(1-y_i)]
+\sum_{i<j}r_{ij}|y_i-y_j|$, and
$\rho_iy_i=\max\{\rho_i,0\}y_i+\max\{-\rho_i,0\}(1-y_i)+\min\{0,\rho_i\}$.
\end{proof}

\begin{theorem}[Certified cut recourse]\label{thm:cr-oracle}
Under~\eqref{eq:cr-class}, for every rational core point $v$, every
rational residual subbox and every rational change of the residual linear
coefficients, one maximum-flow computation returns an exact conditional
minimizer, the exact value $W_{\mathcal C}(v)$, and a certificate: an
endpoint label $y$ and a feasible flow whose value equals
$\operatorname{cap}(S_y)$. The bit complexity is polynomial in the
encoding length of the query, with an absolute exponent independent of
$k$. Given the network, the certificate is checked with
$O(|\mathcal R|+m_{\mathcal R})$ rational operations, where $m_{\mathcal R}$
is the number of nonzero residual couplings. Conditions (R1)--(R2) are
preserved by residual subboxes, by fixing residual coordinates, by added
linear terms and by fixing the core.
\end{theorem}
\begin{proof}
By Lemma~\ref{lem:cr-endpoint} the conditional minimum is attained at an
endpoint label, and by~\eqref{eq:cr-cutid} minimizing over labels is a
minimum $\mathsf s$--$\mathsf t$ cut problem with nonnegative capacities.
Any feasible flow has value at most the capacity of any cut (weak
duality), so a flow whose value equals $\operatorname{cap}(S_y)$ proves
that $S_y$ is a minimum cut; checking it needs capacity bounds,
conservation at each residual node and one comparison. A maximum flow and
such a cut exist and are found, for instance, by the Edmonds--Karp
algorithm with $O(|V||A|^2)$ arithmetic operations, where
$|V|=|\mathcal R|+2$ and $|A|=O(|\mathcal R|+m_{\mathcal R})$. After
multiplying all capacities by a common denominator (whose bit length is
polynomial in the query), capacities are integers, every augmentation is
by an integer, and every flow value is an integer bounded by the total
capacity. Thus all intermediate numbers have polynomial bit length. The
coefficients $\varphi_0,\mu_i,\omega_{ij}$ are obtained by exact
rational evaluation. The stability statement holds because (R1)--(R2)
involve only $a_i$ and $c_{ij}$ for residual $i,j$, which these operations
do not change; a fixed coordinate has $d_i=0$.
\end{proof}

\begin{remark}[Scope and prior work]\label{rem:cr-cut-prior}
Each ingredient is classical: endpoint optimality of coordinatewise
concave functions; the representation of quadratic pseudo-Boolean
functions with nonpositive couplings by cuts
\cite{Hammer1965,PicardRatliff1975,KolmogorovZabih2004}; and the resulting
polynomial solvability of box QPs that are submodular with nonpositive
diagonal, as recorded in \cite{BurerNatarajanWillemsen2026v3}. The role of
Theorem~\ref{thm:cr-oracle} here is a box-stable, certified conditional
oracle with an arbitrary core. Separate concavity (R1) is essential:
continuous submodular box QP without it is not known to be polynomial,
and its natural SDP relaxation is exact only up to dimension three
\cite{BurerNatarajanWillemsen2026v3}. Del Pia and Khajavirad also separate
nonpositive-diagonal variables, which can be taken binary, from the others
and prove polynomial classes with few positive-diagonal variables; their
conditions bound the treewidth and interfaces of the binary part rather
than its signs \cite{DelPiaKhajavirad2026}.
\end{remark}

\subsection{Search over a continuous core}\label{sec:cr-search}

Let the core domain be $[0,1]^k$ ($k\ge1$, continuous) and the residual
domain $X_{\mathcal R}$ any compact mixed box. Put
$V(v)=W_{\mathcal C}(v)=\min_{z\in X_{\mathcal R}}F(v,z)$ and let $L\ge0$ be an
upper coordinate curvature in the core coordinates (for a quadratic,
$L\ge\max_{i\in\mathcal C}2a_i$). Assume an \emph{exact oracle}: for every
dyadic $v\in[0,1]^k$ it returns $V(v)$ and a completion $z(v)\in X_{\mathcal R}$
with $F(v,z(v))=V(v)$, with a certificate. Theorem~\ref{thm:cr-oracle}
supplies one under~\eqref{eq:cr-class}.

\paragraph{Core search.}
Level $j$ uses dyadic cells of side $h_j=2^{-j}$ and
$e_j=kLh_j^2/8$. Start with $\mathcal Q_0=\{[0,1]^k\}$ and $U=+\infty$.
At level $j$: query $V$ at all corners of the cells of $\mathcal Q_j$ that
were not queried before; set $U$ to the least value of all completions
found so far; let
$\lambda_j=\min_{C\in\mathcal Q_j}(\min_{\text{corners of }C}V-e_j)$;
retain $\mathcal K_j=\{C\in\mathcal Q_j:\min_{\text{corners}}V-e_j\le U\}$;
and let $\mathcal Q_{j+1}$ be the $2^k$ dyadic children of each retained
cell. Stop when $U-\min\{\lambda_j,U\}\le\varepsilon$ or at a level limit.

\begin{theorem}[Core search]\label{thm:cr-search}
For every level $j$:
\begin{enumerate}
\item[(a)] every cell of $\mathcal Q_j$ containing $x^*_{\mathcal C}$ for a
global optimizer $x^*$ is retained; in particular $\mathcal K_j\ne\emptyset$;
\item[(b)] $\lambda_j\le\OPT\le U\le\lambda_j+e_j$;
\item[(c)] every retained cell has a corner $v$ with $V(v)\le\OPT+2e_j$;
\item[(d)] every $x\in X$ satisfies $F(x)\ge\min\{\lambda_j,U\}$, and this
can be verified from the list of cells, the certified queried values and
the retention decisions, without any optimization.
\end{enumerate}
\end{theorem}
\begin{proof}
Core coordinates are continuous, so every cell of $\mathcal Q_j$ has
$e_{\mathcal C}(C)=e_j$. (a) The unit cube contains $x^*_{\mathcal C}$. If
$C\in\mathcal Q_j$ contains it, Lemma~\ref{lem:cr-cell} gives
$\min_{\text{corners}}V-e_j\le\OPT\le U$, so $C$ is retained and the child
containing $x^*_{\mathcal C}$ belongs to $\mathcal Q_{j+1}$.
(b)--(c) This is Theorem~\ref{thm:cr-filter} with $\ell=V$ and $\eta=0$;
its hypothesis on an optimizer cell holds by (a).
(d) Children of retained cells cover them, so the cells of $\mathcal Q_j$
together with the cells discarded at levels $i<j$ cover $[0,1]^k$. If
$x_{\mathcal C}$ lies in a cell of $\mathcal Q_j$, then
$F(x)\ge V(x_{\mathcal C})\ge\lambda_j$ by Lemma~\ref{lem:cr-cell}. If it lies in
a cell discarded at level $i$, then
$F(x)\ge\min_{\text{corners}}V-e_i>U^{(i)}\ge U$, where $U^{(i)}$ is the
incumbent at level $i$. Each step uses only recorded values.
\end{proof}

For any $\varepsilon>0$ the search therefore stops by the first level with
$e_j\le\varepsilon$. Its cost is controlled by the number of near-optimal
grid points. Let $N_j$ be the number of points
$v\in h_j\Z^k\cap[0,1]^k$ with $V(v)\le\OPT+2e_j$.

\begin{lemma}[Charging queries]\label{lem:cr-charge}
Level $0$ makes $2^k$ queries and level $j+1$ makes at most $8^kN_j$.
\end{lemma}
\begin{proof}
By Theorem~\ref{thm:cr-search}(c), every retained level-$j$ cell has a
corner counted in $N_j$, and a grid point is a corner of at most $2^k$
cells; so $|\mathcal K_j|\le2^kN_j$. Each retained cell has $2^k$
children with $2^k$ corners each.
\end{proof}

\begin{proposition}[Deterministic count under core growth]\label{prop:cr-growth}
Suppose $V(v)-\OPT\ge g\norm{v-v^*}^2$ for all $v\in[0,1]^k$, some $v^*$
and some $g>0$; this holds with the same $g$ if $F$ has point growth $g$.
Let $\kappa=\max\{1,L/g\}$. Then levels $0,\ldots,J$ make at most
\[
2^k+8^kJ\bigl(1+\sqrt{k\kappa}\bigr)^k
\]
queries. The constant $g$ is not used by the algorithm.
\end{proposition}
\begin{proof}
If $F(x)-\OPT\ge g\norm{x-x^*}^2$, then minimizing over the residual gives
$V(v)-\OPT\ge g\norm{v-x^*_{\mathcal C}}^2$. A point counted in $N_j$
satisfies $g\norm{v-v^*}^2\le2e_j=kLh_j^2/4$, so each coordinate lies
within $\frac{h_j}2\sqrt{k\kappa}$ of $v^*_i$. An interval of length
$h_j\sqrt{k\kappa}$ contains at most $1+\sqrt{k\kappa}$ points of
$h_j\Z$. Hence $N_j\le(1+\sqrt{k\kappa})^k$; apply
Lemma~\ref{lem:cr-charge} and sum over $j<J$.
\end{proof}

\begin{example}[Flat directions]\label{ex:cr-flat}
Without growth the count can be exponential in the number of accuracy
bits. Take $k=2r$, no residual, and $V(v)=\sum_{t=1}^r(v_{2t-1}-v_{2t})^2$,
so $L=2$ and $\OPT=0$. Every product of diagonal cells
$\prod_t\bigl([a_th_j,(a_t+1)h_j]\times[a_th_j,(a_t+1)h_j]\bigr)$ has a
corner with $V=0$, hence is retained. So level $j$ retains at least
$2^{jr}=2^{jk/2}$ cells, and gap $\varepsilon$ costs
$\Omega((L/\varepsilon)^{k/4})$ queries.
\end{example}

Random linear perturbations of the core remove this effect in
expectation. For $\sigma>0$ and an integer $M\ge2$ let $\Gamma_M(\sigma)$
be the uniform distribution on the $M$ points
$-\sigma+2\sigma t/(M-1)$, $t=0,\ldots,M-1$.

\begin{lemma}[Finite uniform law]\label{lem:cr-law}
If $\gamma\sim\Gamma_M(\sigma)$ and $I$ is an interval of length
$\ell\ge0$, then $\Pr[\gamma\in I]\le\ell/(2\sigma)+1/M$.
\end{lemma}
\begin{proof}
Consecutive atoms are $2\sigma/(M-1)$ apart, so $I$ contains at most
$\ell(M-1)/(2\sigma)+1$ atoms. Divide by $M$ and use $(M-1)/M\le1$.
\end{proof}

\begin{theorem}[Smoothed query count]\label{thm:cr-smoothed}
Let $F_\gamma(x)=F_0(x)+\sum_{i\in\mathcal C}\gamma_ix_i$, where $F_0$ is
fixed (it may contain residual data drawn independently of $\gamma$)
and the $\gamma_i$ are independent with law $\Gamma_M(\sigma)$. Let $J\ge1$
and $M\ge2^{J-1}$. The core search is correct for every draw, and the
expected number of queries made by levels $0,\ldots,J$ is at most
\[
2^k+8^kJH_k,\qquad
H_k=\Bigl[3+\Bigl(1+\frac k2\Bigr)\frac L{2\sigma}\Bigr]^k .
\]
\end{theorem}
\begin{proof}
Correctness is Theorem~\ref{thm:cr-search}. Write
$V_\gamma(v)=V_0(v)+\gamma^{\mathsf T}v$ with
$V_0(v)=\min_zF_0(v,z)$; $V_0$ does not depend on $\gamma$, and by
Lemma~\ref{lem:cr-semiconcave} each $t\mapsto V_0(v+t\mathbf e_i)-Lt^2/2$ is
concave. Fix $j\le J-1$, $h=h_j$, a grid point $v$ and a coordinate $i$
with $v_i\in(0,1)$, so $v\pm h\mathbf e_i\in[0,1]^k$. If
$V_\gamma(v)\le\OPT_\gamma+2e_j$, then
$V_\gamma(v)\le V_\gamma(v\pm h\mathbf e_i)+2e_j$, that is,
\[
\frac{V_0(v)-V_0(v+h\mathbf e_i)-2e_j}{h}\ \le\ \gamma_i\ \le\
\frac{V_0(v-h\mathbf e_i)-V_0(v)+2e_j}{h}.
\]
This interval $I_i(v)$ depends on $V_0$, $v$ and $i$ only. Its length is
$[V_0(v+h\mathbf e_i)+V_0(v-h\mathbf e_i)-2V_0(v)+4e_j]/h\le Lh+4e_j/h
=Lh(1+k/2)$, by concavity of $t\mapsto V_0(v+t\mathbf e_i)-Lt^2/2$ at
$t=-h,0,h$. By independence and Lemma~\ref{lem:cr-law},
\[
\Pr\bigl[V_\gamma(v)\le\OPT_\gamma+2e_j\bigr]\le
\prod_{i:\,v_i\in(0,1)}\Bigl[\frac{Lh(1+k/2)}{2\sigma}+\frac1M\Bigr].
\]
Summing over the $(2^j+1)^k$ grid points coordinate by coordinate, with
two boundary values (factor at most $1$) and $2^j-1$ interior values,
\[
\E N_j\le\Bigl[2+(2^j-1)\Bigl(\frac{Lh_j(1+k/2)}{2\sigma}+\frac1M\Bigr)\Bigr]^k
\le H_k ,
\]
since $(2^j-1)h_j\le1$ and $2^j-1<M$. Lemma~\ref{lem:cr-charge} and
summation over $j=0,\ldots,J-1$ complete the proof.
\end{proof}

\begin{remark}[Meaning of the smoothed count]\label{rem:cr-smoothed}
Correctness never depends on the draw, and the instance being solved is
the perturbed one. The bound is linear in the number $J$ of levels,
whereas Example~\ref{ex:cr-flat} shows that the deterministic count can be
exponential in $J$. The residual dimension and graph enter only the
polynomial cost of each query. Each query has encoding length polynomial
in $I+J+\log M$, so with the cut oracle the expected bit work is
$8^kH_k\poly(I+\log(1/\varepsilon)+\log M)$ with an absolute exponent.
The bound is not a device for approximating the unperturbed problem: the
perturbation can change values by $k\sigma$, so accuracy $\delta$ would
require $\sigma\le\delta/k$, and then $H_k$ grows like $\delta^{-k}$,
worse than the deterministic worst case $(kL/\delta)^{k/2}$. With an
approximate oracle satisfying~\eqref{eq:cr-oracle} with $\eta_j\le ae_j$,
the same proof (applied to the true values $V_\gamma$) gives the bound
with $1+k/2$ replaced by $1+(1+a)k/2$. The noise-confinement argument is
in the spirit of isolation lemmas in smoothed analysis
\cite{BeierVocking2004,RoglinVocking2007}.
\end{remark}

For one core coordinate with sign-consistent couplings the
perturbation is unnecessary.

\begin{proposition}[One core coordinate with consistent signs]\label{prop:cr-k1}
Let $k=1$ with core coordinate $v\in[0,1]$, assume~\eqref{eq:cr-class}
with orientation $o$, and assume $c_{vi}o_i\le0$ for every $i\in\mathcal R$
(if instead $c_{vi}o_i\ge0$ for every $i$, substitute $v\mapsto1-v$).
Then an exact global minimizer and $\OPT$ are computed with at most
$2|\mathcal R|+1$ minimum-cut computations and at most $|\mathcal R|+1$
minimizations of univariate quadratics on $[0,1]$, in polynomial time.
\end{proposition}
\begin{proof}
For $S\subseteq\mathcal R$ let $x(S)=\alpha+d\circ\mathbf 1_S$. Collecting
terms as in Proposition~\ref{prop:cr-cut} gives
$F(v,x(S))=q(v)+E_0(S)+v\beta(S)$, where $q$ is a univariate quadratic,
$E_0(S)=\sum_{i\in S}\mu^0_i+\sum_{i<j;\,i,j\in S}\omega_{ij}$ with
$\omega_{ij}\le0$, and $\beta(S)=\sum_{i\in S}\beta_i$ with
$\beta_i=c_{vi}d_i=c_{vi}o_i(u_i-l_i)\le0$. Put
$E_v(S)=E_0(S)+v\beta(S)$ and $\Phi(v)=\min_SE_v(S)$. By
Lemma~\ref{lem:cr-endpoint}, $V=q+\Phi$. The function $E_0$ is submodular
(each $\omega_{ij}y_iy_j$ with $\omega_{ij}\le0$ is) and $\beta$ is
modular, so the minimizers of $E_v$ are closed under union and
intersection; let $S(v)$ be the least one.

\emph{Monotonicity.} Let $v<v'$, $S\in\arg\min E_v$, $T\in\arg\min E_{v'}$.
Submodularity and modularity give
\[
E_v(S\cap T)+E_{v'}(S\cup T)\le E_v(S)+E_{v'}(T)+(v'-v)\beta(S\setminus T)
\le E_v(S)+E_{v'}(T).
\]
Both terms on the left are at least their counterparts on the right, so
$S\cap T\in\arg\min E_v$ and $S\cup T\in\arg\min E_{v'}$. With $S=S(v)$
and $T=S(v')$, minimality gives $S(v)\cap S(v')=S(v)$, that is,
$S(v)\subseteq S(v')$. Hence $\{S(v):v\in[0,1]\}$ is a chain with at most
$|\mathcal R|+1$ members. The source side of a minimum cut reachable from
$\mathsf s$ in the residual network of a maximum flow is the least
minimum cut, so one flow computation returns $S(v)$.

\emph{Recursion.} Write $\ell_S(v)=E_0(S)+v\beta(S)$; then
$\Phi=\min_S\ell_S$ is concave and every $\ell_S\ge\Phi$. Compute
$P=S(0)$ and $Q=S(1)$ and apply the following procedure
$\mathrm{Ex}(0,1,P,Q)$. Given $a<b$ and $P=S(a)\subseteq Q=S(b)$,
$\mathrm{Ex}(a,b,P,Q)$:
\begin{enumerate}
\item stops if $\ell_P$ and $\ell_Q$ do not meet at a single point
$v_m\in(a,b)$;
\item otherwise computes $S_m=S(v_m)$ and stops if
$\Phi(v_m)=\ell_P(v_m)$;
\item otherwise records $S_m$ and calls $\mathrm{Ex}(a,v_m,P,S_m)$ and
$\mathrm{Ex}(v_m,b,S_m,Q)$.
\end{enumerate}

At a stop, $\Phi=\min\{\ell_P,\ell_Q\}$ on $[a,b]$. Indeed $\Phi$ is
concave, $\Phi\le\ell_P,\ell_Q$, $\Phi(a)=\ell_P(a)$ and
$\Phi(b)=\ell_Q(b)$. A concave function lying below a line and touching
it at both ends of an interval equals the line there, because it lies
above the chord. Since $P\subseteq Q$ and $\beta\le0$, the slope of
$\ell_P$ is at least that of $\ell_Q$. In the second stopping case apply
the chord fact on $[a,v_m]$ to $\ell_P$ and on $[v_m,b]$ to $\ell_Q$; by
the slope comparison $\min\{\ell_P,\ell_Q\}$ is $\ell_P$ left of $v_m$ and
$\ell_Q$ right of it. In the first case the slope comparison together with
$\ell_P(a)\le\ell_Q(a)$ and $\ell_Q(b)\le\ell_P(b)$ shows that one line,
say $\ell_Q$, lies below the other on all of $[a,b]$ (the other case is
symmetric); then $\ell_Q(a)\le\ell_P(a)=\Phi(a)\le\ell_Q(a)$, so $\ell_Q$
touches $\Phi$ at both ends and $\Phi=\ell_Q=\min\{\ell_P,\ell_Q\}$ on
$[a,b]$. A recorded $S_m$ satisfies
$P\subseteq S_m\subseteq Q$ and differs from both, since
$\ell_{S_m}(v_m)<\ell_P(v_m)=\ell_Q(v_m)$; sets recorded in different
branches lie in disjoint parts of the chain. So at most $|\mathcal R|-1$
sets are recorded, the recursion tree has at most $|\mathcal R|$ leaves,
and with the two initial computations there are at most
$2|\mathcal R|+1$ flow computations. Every $v_m$ is the intersection of two
lines with rational coefficients of polynomial bit length.

Let $\mathcal S$ be $P$, $Q$ and the recorded sets. Then
$\Phi=\min_{S\in\mathcal S}\ell_S$ on $[0,1]$, so
$\OPT=\min_{S\in\mathcal S}\min_{v\in[0,1]}(q(v)+\ell_S(v))$, a minimum of
at most $|\mathcal R|+1$ univariate quadratics on $[0,1]$, each solved
exactly.
\end{proof}

Without the sign condition the least minimizers need not be nested, and
we do not know whether $k=1$ admits a deterministic polynomial algorithm.

\subsection{Exact output}\label{sec:cr-exact}

\begin{lemma}[Uniform height over residual labels]\label{lem:cr-height}
Let $F$ be a rational quadratic with continuous core $[0,1]^k$ and residual
satisfying (R1), with integer residual bounds rounded inward. Let
$\Lambda_c$ be the least common denominator of all coefficients of $F$,
$\Lambda_e$ that of all residual bounds, $T_0=\Lambda_c\Lambda_e^2$,
$H_0=\max\{1,\max_{i,j\in\mathcal C}|T_0\,\partial_{ij}F|\}$ and
\[
D_0=T_0\bigl(k!\,H_0^k\bigr)^2 .
\]
For every residual endpoint label $z$, the value
$\Phi(z)=\min_{v\in[0,1]^k}F(v,z)$ is a rational with reduced denominator
at most $D_0$, and so is $\OPT$.
\end{lemma}
\begin{proof}
Fix $z$ and put $\phi(v)=F(v,z)$. Every monomial of $T_0\phi$ has an
integer coefficient: the coefficients of $v_iv_j$ and $v_i^2$ are
coefficients of $F$; the coefficient of $v_i$ is $b_i+\sum_jc_{ij}z_j$;
and the constant collects $c_0$, $a_jz_j^2$, $b_jz_j$ and $c_{ij}z_iz_j$
for residual $i,j$. Each has denominator dividing $\Lambda_c\Lambda_e^2$.
The matrix $A=T_0\nabla^2_{\mathcal C\mathcal C}F$ is integral with entries
bounded by $H_0$ in absolute value.

Among minimizers of $\phi$ on $[0,1]^k$ choose one, $v$, with the fewest
coordinates in $(0,1)$; let $S$ be that set. The other coordinates are
$0$ or $1$. If $S=\emptyset$, then $T_0\phi(v)\in\Z$. Otherwise
$\nabla_S\phi(v)=0$ and $A_{SS}\succeq0$ by first- and second-order
optimality on the face. If $A_{SS}$ were singular, a nonzero $w$
supported on $S$ with $A_{SS}w_S=0$ would give $\phi(v+tw)=\phi(v)$ for all
$t$, and moving $t$ until a coordinate of $S$ reaches $0$ or $1$ would
produce a minimizer with fewer interior coordinates. Hence
$\delta=|\det A_{SS}|$ is a positive integer, and by the Leibniz expansion
$\delta\le|S|!\,H_0^{|S|}\le k!\,H_0^k$. The stationarity system
$A_{SS}v_S=-T_0\nabla_S\phi(0_S,v_{\mathcal C\setminus S})$ has an integral
right-hand side, so $\delta v$ is integral by Cramer's rule. Evaluating
the integer-coefficient quadratic $T_0\phi$ at $v=(\delta v)/\delta$ gives
an element of $\delta^{-2}\Z$. Thus $\Phi(z)\in(T_0\delta^2)^{-1}\Z$.
By Lemma~\ref{lem:cr-endpoint}, some global optimizer has an endpoint
label $z^*$, and $\OPT=\Phi(z^*)$.
\end{proof}

\begin{remark}
Any common multiple $S_0$ of all denominators can replace $\Lambda_c$ and
$\Lambda_e$, for example $T_0=S_0^3$. A cube, or the mixed form
$\Lambda_c\Lambda_e^2$, is needed in general: the term $\frac12x_ix_j$ at
residual endpoints $x_i=x_j=\frac12$ equals $\frac18$, while the least
common denominator of the data is $2$.
\end{remark}

\begin{theorem}[Exact output with a cut residual]\label{thm:cr-exact}
Assume~\eqref{eq:cr-class} with continuous core $[0,1]^k$ and rational data.
Run the core search until $U-\min\{\lambda_j,U\}<D_0^{-2}$, which happens
by the first level with $e_j<D_0^{-2}$. Let $z_U$ be the residual label
of the incumbent completion. Minimize $\phi(v)=F(v,z_U)$ exactly over
$[0,1]^k$ by enumerating the $3^k$ faces of the cube, solving the
stationarity system on each face whose free Hessian block is nonsingular,
and keeping feasible solutions. The best candidate $\hat v$ satisfies
$F(\hat v,z_U)=\OPT$, so $(\hat v,z_U)$ is a global optimizer. A checker
needs only the search trace, $D_0$, feasibility and exact value of the
output, its reduced denominator ($\le D_0$), and the inequality
$F(\hat v,z_U)-\min\{\lambda_j,U\}<D_0^{-2}$.
If in addition $V$ has core growth with constant $g$ as in
Proposition~\ref{prop:cr-growth}, the total bit complexity is
$f(k,\kappa)(I+1)^{C_1}$ with an absolute constant $C_1$, and $g$ is not
an input.
\end{theorem}
\begin{proof}
The proof of Lemma~\ref{lem:cr-height} shows that a minimizer of $\phi$
lies on some face with a nonsingular free block, so the enumeration
contains it; all candidates are feasible, so the best one attains
$\Phi(z_U)$. If $(v_U,z_U)$ is the incumbent completion, then
$\OPT\le\Phi(z_U)\le F(v_U,z_U)=U$ and, by Theorem~\ref{thm:cr-search}(d),
$\min\{\lambda_j,U\}\le\OPT$. So $0\le\Phi(z_U)-\OPT<D_0^{-2}$. Both
numbers have reduced denominators at most $D_0$ by
Lemma~\ref{lem:cr-height}, and distinct such rationals differ by at least
$D_0^{-2}$; hence they are equal. The checker's test proves optimality of
any feasible point with these properties by the same separation argument.

For the complexity, $\log T_0$ and $\log H_0$ are $O(I)$ and $k\le I$, so
$\log D_0=O(I^2)$ and the number of levels is
$J=O(\log D_0+\log(kL)+1)=\poly(I)$. Under core growth,
Proposition~\ref{prop:cr-growth} bounds the queries by
$2^k+8^kJ(1+\sqrt{k\kappa})^k$. Each query is a cut problem whose data have
bit length polynomial in $I+J$, hence $\poly(I)$ by
Theorem~\ref{thm:cr-oracle}. Face enumeration costs $3^k\poly(I)$. All
factors depending only on $k$ and $\kappa$ form $f$.
\end{proof}

Without core growth the exact mode still terminates, but its number of
cells can be $2^{kJ}$ with $J=\Theta(\log D_0)$. The stopping depth must
not be combined with Theorem~\ref{thm:cr-smoothed} by choosing
$M\ge2^{J}$: sampled coefficients enter $D_0$, and for a fixed noise scale
the atom denominators grow with $M$, so this condition generally cannot
hold.

\subsection{Mixed concave--convex residuals}\label{sec:cr-mixed}

Consider a rational residual problem
$\min\{q(z)=\frac12z^{\mathsf T}Az+b^{\mathsf T}z+c : z\in\prod_{i\in\mathcal R}[l_i,u_i]\}$,
for example the residual after the core has been fixed. Assume an
orientation $o$ with $o_io_jA_{ij}\le0$ for $i\ne j$, and a partition
$\mathcal R=\mathcal D\cup\mathcal P$ with $A_{ii}\le0$ for $i\in\mathcal D$ and
$A_{\mathcal P\mathcal P}\succeq0$. Coordinates in $\mathcal P$ are
continuous; coordinates in $\mathcal D$ may be integer with inward-rounded
bounds. Define $\alpha_i,d_i$ as before and, for
$S\subseteq\mathcal D$, the label $\zeta(S)_i=\alpha_i+d_i[i\in S]$ and
\[
g(S)=\min\Bigl\{q(\zeta(S),z_{\mathcal P}) :
z_{\mathcal P}\in\prod_{i\in\mathcal P}[l_i,u_i]\Bigr\}.
\]
With $\mathcal P=\emptyset$ this is the cut class; with
$\mathcal D=\emptyset$ it is convex.

\begin{proposition}[Submodular conditional value]\label{prop:cr-submod}
(a) The minimum of $q$ over the continuous box and over the mixed domain
equals $\min_{S\subseteq\mathcal D}g(S)$.
(b) Each $g(S)$ is the value of a rational convex box QP.
(c) For all $S,T\subseteq\mathcal D$,
$g(S)+g(T)\ge g(S\cap T)+g(S\cup T)$.
\end{proposition}
\begin{proof}
(a) As in Lemma~\ref{lem:cr-endpoint}, move each coordinate of
$\mathcal D$ to an endpoint without increasing $q$. (b) After fixing
$z_{\mathcal D}$, the objective in $z_{\mathcal P}$ has Hessian
$A_{\mathcal P\mathcal P}\succeq0$. (c) Substitute out degenerate intervals
and put $z_i=\alpha_i+d_it_i$, $t\in[0,1]^{\mathcal R}$; the transformed objective
$\tilde q$ has off-diagonal Hessian entries
$d_id_jA_{ij}=o_io_j(u_i-l_i)(u_j-l_j)A_{ij}\le0$. For $s,t\in[0,1]^{\mathcal R}$,
$\tilde q(s)+\tilde q(t)\ge\tilde q(s\wedge t)+\tilde q(s\vee t)$: unary terms
cancel, and for a pair $i,j$ the two sides agree if $s_i-t_i$ and
$s_j-t_j$ have the same sign, while otherwise
$s_is_j+t_it_j-(s\wedge t)_i(s\wedge t)_j-(s\vee t)_i(s\vee t)_j
=(s_i-t_i)(s_j-t_j)\le0$, which multiplied by a nonpositive coefficient is
nonnegative. Let $t^S,t^T$ attain $g(S),g(T)$ in transformed
coordinates. Then $(\mathbf 1_S,t^S)\wedge(\mathbf 1_T,t^T)
=(\mathbf 1_{S\cap T},t^S\wedge t^T)$ and similarly for $\vee$ with
$S\cup T$; both completions are feasible. Hence
$g(S)+g(T)\ge\tilde q(\mathbf 1_{S\cap T},t^S\wedge t^T)
+\tilde q(\mathbf 1_{S\cup T},t^S\vee t^T)\ge g(S\cap T)+g(S\cup T)$.
\end{proof}

Positive semidefiniteness is not used in (c); it makes each $g(S)$
computable exactly in polynomial time \cite{KozlovTarasovKhachiyan1980}.
Partial minimization preserving submodularity is classical
\cite{Topkis1978}, and the same composition with submodular minimization
appears in \cite{BuntonTabuada2022,GomezHan2025}.

Let $m=|\mathcal D|$ and $h(S)=g(S)-g(\emptyset)$. For an ordering $\pi$ of
$\mathcal D$ with prefixes $S_0=\emptyset,S_1,\ldots,S_m=\mathcal D$, the
greedy vector is $b^\pi_{\pi(t)}=h(S_t)-h(S_{t-1})$.

\begin{proposition}[Greedy-base certificate]\label{prop:cr-greedy}
(a) For every ordering $\pi$ and every $T\subseteq\mathcal D$,
$b^\pi(T)\le h(T)$ and $b^\pi(\mathcal D)=h(\mathcal D)$.
(b) If $w=\sum_\pi\lambda_\pi b^\pi$ is a convex combination, then
$\min q\ge g(\emptyset)+\sum_{i\in\mathcal D}\min\{0,w_i\}$.
(c) Some convex combination of at most $m+1$ greedy vectors attains
equality \cite{Edmonds1970}, and such a combination, a minimizing label
and exact convex-QP witnesses can be computed in time polynomial in the
input length \cite{GrotschelLovaszSchrijver1988}.
(d) Given orderings, weights, the $O(m^2)$ chain values with their
minimizers, and a label $S^*$ with its completion, a checker verifies the
bound in (b) and its attainment by $S^*$ with polynomially many rational
operations and no optimization. It checks a convex chain value by
feasibility, exact evaluation and the sign conditions
$\partial_iq=0$ at interior, $\ge0$ at lower and $\le0$ at upper bounds
of $\mathcal P$, which suffice by convexity once
$A_{\mathcal P\mathcal P}\succeq0$ has been verified exactly.
\end{proposition}
\begin{proof}
(a) List $T=\{t_1,\ldots,t_r\}$ in $\pi$-order and let $S^{(s)}$ be the
$\pi$-prefix ending just before $t_s$. Since
$\{t_1,\ldots,t_{s-1}\}\subseteq S^{(s)}$, submodularity gives
$b^\pi_{t_s}=h(S^{(s)}\cup\{t_s\})-h(S^{(s)})
\le h(\{t_1,\ldots,t_s\})-h(\{t_1,\ldots,t_{s-1}\})$. Summing telescopes
to $b^\pi(T)\le h(T)$; the full prefix sum is $h(\mathcal D)$.
(b) By (a) and convexity, $w(S)\le h(S)$ for every $S$, so
$g(S)=g(\emptyset)+h(S)\ge g(\emptyset)+w(S)\ge g(\emptyset)+\sum_i\min\{0,w_i\}$;
conclude with Proposition~\ref{prop:cr-submod}(a).
(c) Edmonds' theorem gives
$\min_Sh(S)=\max\{\sum_i\min\{0,w_i\}: w\in B(h)\}$, where $B(h)$ is the
base polytope; its vertices are the greedy vectors, and the maximum of
this concave piecewise-linear function is attained at a combination of at
most $m+1$ of them. For computation, consider the linear program
$\max\{\beta: 0\le y\le\mathbf 1,\ \beta+(b^\pi)^{\mathsf T}y\le0\ \forall\pi\}$,
whose dual is $\min\{\sum_ir_i : \sum_\pi\lambda_\pi b^\pi+r\ge0,\
\sum_\pi\lambda_\pi=1,\ \lambda,r\ge0\}$. Sorting $y$ decreasingly and
forming its greedy vector maximizes $(b^\pi)^{\mathsf T}y$, so separation
costs $m+1$ convex-QP queries. Each $g(S)$ is the optimal value of a convex
box QP with fixed Hessian, and by the smallest-face argument of
Lemma~\ref{lem:cr-height} its encoding length is polynomial; hence the
polyhedron has polynomial facet complexity, and the ellipsoid method
solves the program and recovers an optimal basic dual solution from the
separating inequalities in polynomial time
\cite[Thm.~6.4.9, Lemma~6.5.15]{GrotschelLovaszSchrijver1988}. A basic dual
solution has at most $m+1$ positive $\lambda_\pi$, and at optimality
$r_i=\max\{0,-w_i\}$. For an optimal $y$ sorted as
$y_{\pi(1)}\ge\cdots\ge y_{\pi(m)}$, with $y_{\pi(m+1)}:=0$,
$(b^\pi)^{\mathsf T}y=\sum_t(y_{\pi(t)}-y_{\pi(t+1)})h(S_t)+(1-y_{\pi(1)})h(\emptyset)$
is a convex combination of prefix values equal to $\min h$, so some prefix
with positive weight is a minimizing label. (d) follows from (a), (b) and
the stated optimality conditions for convex problems on a box.
\end{proof}
